\documentclass[10pt]{article}
\usepackage[a4paper,margin=23mm]{geometry}
\usepackage{amsmath,amssymb,amsthm,booktabs,longtable,hyperref}
\newtheorem{theorem}{Theorem}\newtheorem{lemma}{Lemma}
\title{Literal diagram and construction joins for the $\operatorname{PSL}_3(4)$ catalogue}
\author{Finite source-recovery component}\date{}
\begin{document}\maketitle
\noindent
The fixed source is Gottschalk and Leemans, \emph{Geometries for the group
$\operatorname{PSL}(3,4)$}, Halle Reports on Algebra and Geometry 99--30,
the complete 22-page institutional preprint.  Its PostScript and derived PDF
are retained exactly.  The 25-page journal edition, European Journal of
Combinatorics 24 (2003), 267--291, has not been collated here.  Statements about
the printed diagrams below refer only to the preprint.

\paragraph{Acquired inputs.}
The accepted broad-primitive enumeration exhausts all 95 subgroup anchors in a
fixed permutation group $G$ of order $20160$, using the retained complete
subgroup and maximal-inclusion inventories.  Its full-automorphism fusion has
29 representatives of ranks $2$--$5$.  The separately accepted palette
calculation, together with its exact serialization repair, provides complete
right-coset partitions, incidence rows and actual $G$-actions for these
representatives.  The full-automorphism and GRAPE/nauty completeness premises
remain explicit acquired inputs.  They are not proved by checking positive
permutation maps.  The previously accepted projective construction component
provides two rank-five models and eight literal truncation routes.

\begin{theorem}
The retained literal calculations establish the following statements.
\begin{enumerate}
\item The numerical node data, all rank-two residues, local action kernels
and local $2$-transitivity flags are determined for each of the 29 acquired
representatives.  Within this complete catalogue, every preprint diagram has
one representative matching its recorded invariants, subject to the two
discrepancies in Lemma~\ref{errors}.
\item All 52 printed circular, Petersen and affine-plane bond labels have
explicit canonical incidence maps, including their displayed polarity.
\item The accepted constructions labelled 5.1 and 5.2 are explicitly isomorphic
to catalogue representatives 26 and 28, respectively.  The eight acquired
truncation routes also have complete object and incidence maps.
\end{enumerate}
This is not an assertion that diagrams determine geometries in general.
Abstract names of Borel groups, parabolics, kernels, automorphism groups and
correlation groups require their own structural identifications.
\end{theorem}

\begin{proof}
Choose the identity-coset chamber.  For a set $J$ of remaining types, intersect
the neighbour sets of all complementary chamber objects.  These are exactly
the objects of the base residue of type $J$.  Acquired flag transitivity
transports this calculation to every flag of the same type.  A panel order is
its cardinality minus one.  In its actual stabilizer, the stabilizer of one
panel object is tested for transitivity on the other objects.  No local
$2$-transitivity assumption is imposed.

For every pair of types, all incidence and nonincidence rows of the resulting
bipartite graph are retained.  Breadth-first distances give eccentricities, compared at both ordered
endpoints of each printed numerical $a/g/b$ label, and the minimum cycle length; half that length is the gonality.  A residue is a
digon precisely when it is complete bipartite.  At each base object, every
element of its parabolic is tested on every object of its residue.  The
elements fixing all these objects are exactly the local action kernel.
Only its order, not its abstract structure, is used in the numerical join.

For a circular bond, the degree-two objects map bijectively to all unordered
pairs of the opposite object set.  For a Petersen bond, an explicit map sends
the ten points to the two-subsets of a five-set and the fifteen lines to
disjoint pairs of such subsets.  Every incidence row is checked.  For an
affine bond, sixteen points and twenty four-point lines have a unique line
through each point pair and five parallel classes.  Adding one point for each
parallel class and a line at infinity gives a projective completion.  An
explicit 21-point bijection takes all its lines to the coordinate plane over
$\mathbb F_4$.  The recorded point/line maps also verify the source polarity.

Finally, the transcription records actual node positions, not presumed type
numbers.  All permutations compatible with the printed type sizes are checked
against the complete residue data.  The resulting table below records the
unique catalogue index in each case; symmetric node permutations may remain.
These tests compare the complete finite catalogue with the source, rather
than using source counts as acceptance criteria.
\end{proof}

\begin{center}\small
\begin{tabular}{@{}lrrrrrrrrrrr@{}}\toprule
Source rank 2 &2.1&2.2&2.3&2.4&2.5&2.6&2.7\\
Catalogue &1&4&2&5&3&6&7\\\midrule
Source rank 3 &3.1&3.2&3.3&3.4&3.5&3.6&3.7\\
Catalogue &9&13&10&11&14&12&8\\\midrule
Source rank 4 &4.1&4.2&4.3&4.4&4.5&4.6&4.7&4.8&4.9&4.10&4.11\\
Catalogue &15&24&19&22&23&18&17&16&21&20&25\\\midrule
Source rank 5 &5.1&5.2&5.3&5.4\\
Catalogue &26&28&29&27\\\bottomrule
\end{tabular}
\end{center}

\paragraph{Construction maps and action conventions.}
The projective component and census use different generator lists.  Their
complete generated permutation groups are verified equal on the same 21
labels.  The accepted coordinate-to-retained map is reused.  If necessary,
an explicit duality is applied: a normal vector labels its perpendicular
projective line, and the induced action on these 21 lines is transported back
to the fixed point labels.  This is checked to permute all $20160$ group
elements bijectively and preserve every edge of the generator Cayley graph,
so it is an automorphism $\alpha$ of the actual group.

For every proposed object $x_i$, the generators of the actual parabolic $H_i$
fix $x_i$ after applying $\alpha$.  Consequently
\[
                  H_i g\longmapsto x_i^{\alpha(g)}
\]
is well defined.  The certificate additionally checks every coset member,
the complete object bijections, all incidence and nonincidence rows, and the
explicit type permutation.  Separate checks of the saved maps verify
3,185 incidence rows and 6,370 actual generator-action rows.

\begin{lemma}\label{errors}
In preprint diagrams 5.1 and 5.2, the printed order $2$ at the unique
210-object node must be $1$ for the stated constructions.
\end{lemma}
\begin{proof}
The literal construction maps send the 210-object node to the unordered
point-pair type.  Its panel has two objects and therefore order $1$, as already
verified in the accepted construction.  The order-$2$ line node has 21 objects
and cannot be confused with this node.

\end{proof}

\paragraph{Scope and reproducibility.}
The scripts retain all 29 literal residue tables, 52 canonical bond maps and
ten construction/truncation maps.  Bounded controls include 25 known graph
distance/girth examples, six complete-graph incidence examples, damaged maps
and missing endpoints.  Two failed preliminary joins are preserved: one
incorrectly required identical generator lists, and one omitted the needed
duality.  Neither negative attempt is a premise.  No graph-isomorphism target
or subgroup census is repeated by this component.

The original preprint has a single horizontal bond at the base of diagram 3.5,
consistent with the literal six-cycle residue.  The first author edition
misread this bond and incorrectly proposed a source correction; that edition
and its passing finite checks are preserved, and this false claim is withdrawn.
The successor additionally compares both ordered endpoint eccentricities
in every numerical $a/g/b$ bond label, not only its middle gonality.
The source has no local-$2$-transitivity marker at 3.5 or 4.11, consistent with
the actual panels.  Earlier provisional transcription mistakes at those two
rows are disclosed and corrected.  The two discrepancies above do not
establish an error in the uncollated journal edition.  Remaining constructions,
the general doubling hypotheses, abstract group-name claims and whole-source
completion are outside this component.  There is no new question, Lean,
measurement or novelty claim.
\end{document}
